
% =====================================================================
\section{T4: every part has a private invariant}
\label{sec:invariant}
% =====================================================================

The residual of \Cref{sec:residual} is an invariant of the whole graph. We now
show the same construction localises: every connected part carries its own
invariant under its own relabellings, computed from its internal boundary
content together with the boundary that joins it to the rest. This per-part
invariant is the standing structure from which the part's surface activity
derives.

\subsection{The internal residual of a part}

\begin{definition}[Private invariant of a part]
\label{def:private}
Let $\Agent=\Gph[U]$ be a part (a connected induced subgraph,
\Cref{def:subgraph}). Its \emph{private invariant} is the pair
\[
\Res_{\mathrm{priv}}(\Agent)\;=\;\bigl(\Res(\Agent),\,b(U)\bigr),
\]
where $\Res(\Agent)$ is the internal residual of the induced subgraph
\textup{(}\Cref{def:residual} applied to $\Agent$\textup{)} and $b(U)=\wt(\cut(U))$
is its boundary cost to the rest of $\Gph$. We treat $\Res_{\mathrm{priv}}$ up to
weighted isomorphism of $\Agent$ as an abstract subgraph.
\end{definition}

\begin{theorem}[Privacy and invariance of the part-invariant]
\label{thm:private}
For every part $\Agent$ of a contact graph:
\begin{enumerate}[label=\textup{(\roman*)},leftmargin=2.2em]
\item \textbf{Positivity.} $\Res(\Agent)\ge\floorc>0$ whenever $\abs{U}\ge2$,
and $b(U)\ge\floorc>0$ always; the private invariant is bounded away from zero in
both components.
\item \textbf{Invariance.} $\Res_{\mathrm{priv}}(\Agent)$ is invariant under
$\mathrm{Aut}(\Agent)$ and under any weighted isomorphism of $\Agent$ that
preserves its boundary cost; it is a property of $\Agent$'s isomorphism class,
not of its labelling.
\item \textbf{Privacy.} Two parts with the same private invariant need not be the
same part of $\Gph$; conversely a part's private invariant is determined by the
part alone, independently of how the rest of $\Gph$ is labelled.
\end{enumerate}
\end{theorem}

\begin{proof}
(i) $\Agent$ is a contact graph in its own right (a finite, connected, weighted
graph with the inherited floor $\floorc$), so \Cref{thm:residual-region}(i) gives
$\Res(\Agent)\ge\floorc$ for $\abs{U}\ge2$; and $b(U)\ge\floorc$ by
\Cref{thm:floor}. (ii) Apply \Cref{thm:residual-invariant} to the contact graph
$\Agent$: $\Res(\Agent)$ is invariant under weighted isomorphisms of $\Agent$;
$b(U)$ is the total weight of the boundary cut, manifestly preserved by any
isomorphism fixing that cut setwise. Hence the pair is invariant on the relevant
class. (iii) The internal residual and boundary cost are computed from $\Agent$
and its incident cut alone; relabelling the complement $\V\setminus U$ changes
neither, establishing independence from the rest's labelling. That distinct parts
may share the invariant is immediate (e.g. two isomorphic clusters in different
positions).
\end{proof}

\subsection{The invariant is the standing structure}

\begin{proposition}[Surface walks derive from the private invariant]
\label{prop:surface-derives}
Fix a part $\Agent=\Gph[U]$. Any $\gate$-admissible walk confined to $\Agent$
\textup{(}\Cref{def:walk,def:gate}\textup{)} traverses only edges of $\Agent$,
whose existence and weights are fixed by $\Agent$'s isomorphism class. Hence the
set of walks available within $\Agent$ is determined by $\Agent$ up to
relabelling, while no individual walk is an invariant of $\Agent$: the walks are
the variable surface, the subgraph \textup{(}with its private invariant\textup{)}
is the standing structure they run on.
\end{proposition}

\begin{proof}
A walk within $\Agent$ uses only edges $uv$ with $u,v\in U$; the availability of
each such step is exactly the presence of that edge in $\Gph[U]$, an
isomorphism-invariant datum. Thus the \emph{space} of available walks transforms
covariantly with $\Agent$ under relabelling (it is carried bijectively by any
isomorphism), so it is determined by $\Agent$'s class. A single walk, by
contrast, is a labelled sequence and is not preserved by a nontrivial
relabelling; hence no individual walk is invariant. The standing/variable split
follows.
\end{proof}

\begin{remark}[Interpretation: the part's ``identity'', on which nothing depends]
\label{rem:character}
\Cref{thm:private} and \Cref{prop:surface-derives} say a part has a conserved
internal structure---its private invariant---distinct from any of the particular
trajectories it may run, and unique to it only up to isomorphism. If one reads a
part as an agent and its walks as the agent's successive states, then the private
invariant is the agent's standing identity: the conserved bottom from which its
changing activity derives, the local counterpart of the global conserved residual
of \Cref{sec:residual}. The reading is offered as orientation; the theorems are
graph-theoretic and used only as such below.
\end{remark}

\begin{corollary}[Private and global invariants are one construction at two
scales]
\label{cor:two-scales}
The global residual $\Res(\Gph)$ \textup{(}\Cref{def:residual}\textup{)} and the
private invariant $\Res_{\mathrm{priv}}(\Agent)$ \textup{(}\Cref{def:private}\textup{)}
are the same functional---minimum inter-part boundary content---evaluated on the
whole graph and on an induced subgraph respectively. Both are positive
\textup{(}T0\textup{)} and invariant under re-encoding \textup{(}T3\textup{)}. The
private invariant is the global construction restricted to a part.
\end{corollary}

\begin{proof}
Immediate by comparing \Cref{def:residual} with \Cref{def:private} and applying
\Cref{thm:residual-region,thm:residual-invariant} to $\Gph$ and to $\Agent$.
\end{proof}
